\section{Radice Quadrata di una Matrice}

In molte applicazioni, in particolare nello studio delle distribuzioni gaussiane
multivariate, compare l’espressione $\sqrt{C}$, dove $C$ è una matrice. In 
questa appendice spieghiamo con precisione cosa significhi ``fare la radice
quadrata'' di una matrice e perché questo concetto è utile.

\begin{definizione}[Radice Quadrata]
    Sia $C$ una matrice quadrata $d\times d$. Una matrice $A$ si dice 
\textit{radice quadrata} di $C$ se soddisfa
\[
A^2 = A A = C.
\]
\end{definizione}


Questa non è in generale univoca: possono esistere molte matrici diverse $A$
tali che $A^2=C$. Tuttavia, quando $C$ è simmetrica e definita positiva, la
situazione diventa molto più semplice.

\bigskip
\noindent Sia $C$ simmetrica definita positiva. Allora:

\begin{itemize}
    \item tutti gli autovalori di $C$ sono reali e positivi,
    \item $C$ è diagonalizzabile ortogonalmente: 
    \[
    C = Q \Lambda Q^\top,
    \]
\end{itemize}
dove $\Lambda=\mathrm{diag}(\lambda_1,\dots,\lambda_d)$ con $\lambda_i>0$, e $Q$ è ortogonale.

\begin{definizione}[Radice Principale]
    In questo caso possiamo definire la \textit{radice quadrata principale} come:
\[
\sqrt{C} := Q\,\sqrt{\Lambda}\,Q^\top,
\]
dove 
\[
\sqrt{\Lambda} = \mathrm{diag}(\sqrt{\lambda_1},\dots,\sqrt{\lambda_d}).
\]
\end{definizione}
\noindent È immediato verificare che:
\[
\sqrt{C}\,\sqrt{C} = C.
\]

\noindent Questa radice quadrata è univoca, simmetrica e definita positiva.

\bigskip
\noindent Se $C$ non è definita positiva o non è simmetrica, possono esistere molte
radici diverse. Ad esempio:
\[
I_2 = 
\begin{pmatrix}
1 & 0 \\ 0 & 1
\end{pmatrix}
\qquad\Rightarrow\qquad
A = 
\begin{pmatrix}
\cos\theta & -\sin\theta \\
\sin\theta & \cos\theta
\end{pmatrix}
\]
soddisfa $A^2 = I_2$ per ogni $\theta$.